%%%%%%%%%%%%%%%%%%%%%%%%%%%%%%%%%%%%%%%%%
% Medium Length Professional CV
% LaTeX Template
% Version 2.0 (8/5/13)
%
% This template has been downloaded from:
% http://www.LaTeXTemplates.com
%
% Original author:
% Rishi Shah 
%
% Important note:
% This template requires the resume.cls file to be in the same directory as the
% .tex file. The resume.cls file provides the resume style used for structuring the
% document.
%
%%%%%%%%%%%%%%%%%%%%%%%%%%%%%%%%%%%%%%%%%

%----------------------------------------------------------------------------------------
%	PACKAGES AND OTHER DOCUMENT CONFIGURATIONS
%----------------------------------------------------------------------------------------

\documentclass{resume} % Use the custom resume.cls style

\usepackage{fontspec}
\setmainfont{Source Sans Pro}

\usepackage{microtype}
\linespread{0.97} %good spacing for resumes

\usepackage[backend=biber,
maxnames=2,
style=authoryear]{biblatex}
\addbibresource{jeherr.bib}

\renewcommand*{\mkbibnamegiven}[1]{%
  \ifitemannotation{highlight}
    {\textbf{#1}}
    {\ifitemannotation{highcorr}
    {\textbf{#1 *}}
    {\ifitemannotation{highequal}
    {\textbf{#1 \textsuperscript{§}}}
    {\ifitemannotation{equal}
    {#1\textbf{\textsuperscript{§}}}
    {#1}}}}}

\renewcommand*{\mkbibnamefamily}[1]{%
  \ifitemannotation{highlight}
    {\textbf{#1}}
    {\ifitemannotation{highcorr}
    {\textbf{#1}}
    {\ifitemannotation{highequal}
    {\textbf{#1}}
    {\ifitemannotation{equal}
    {#1}
    {#1}}}}}



\usepackage{hyperref}
\hypersetup{
    colorlinks=true,
    urlcolor=blue
}

\usepackage[margin=0.6in]{geometry}
\newcommand{\tab}[1]{\hspace{.2667\textwidth}\rlap{#1}}
\newcommand{\itab}[1]{\hspace{0em}\rlap{#1}}
\name{John E. Herr} % Your name

\address{(319) 670-8252 \\ johnherr@gmail.com \\ \url{https://github.com/jeherr}}

\begin{document}


%----------------------------------------------------------------------------------------
%	Summary
%----------------------------------------------------------------------------------------
\begin{rSection}{Summary}

Computational chemist \& machine‑learning scientist with 6+ years building deep‑learning models for molecular property prediction and force‑field development. Developed datasets with over 10 M conformations, authored 12 peer‑reviewed papers (1,700+ citations), and cut CPU cost of QM pipelines by up to 40\%. Proficient in Python, PyTorch, TensorFlow, RDKit, and HPC workflows.

\end{rSection}

%----------------------------------------------------------------------------------------
%	Technical Skills
%----------------------------------------------------------------------------------------

\begin{rSection}{Technical Skills}

\begin{tabular}{ @{} l l @{} }
\textbf{Cheminformatics:} & reaction prediction, large-scale data curation \\
\textbf{Simulation Software:} & Q-Chem, ORCA, Psi4, PySCF, CP2K, VASP, Quantum ESPRESSO \\
\textbf{Machine Learning:} & Graph networks, transformers, autoencoders, GANs \\
\textbf{Programming:} & Python, PyTorch, TensorFlow, RDKit \\
\end{tabular}

\end{rSection}


%----------------------------------------------------------------------------------------
%	Education
%----------------------------------------------------------------------------------------

\begin{rSection}{Education}

{\bf University of Notre Dame du Lac} \hfill {\em August 2015 - May 2020}
\newline Ph.D. Theoretical and Computational Chemistry

{\bf University of Iowa}   \hfill {\em August 2011 - May 2014} 
\newline B.S. Chemistry and Mathematics

% {\bf Southeastern Community College} \hfill {\em August 2007 - July 2010}
% \newline Associate of Arts

\end{rSection}

%----------------------------------------------------------------------------------------
%	Research Experience
%----------------------------------------------------------------------------------------

% \begin{rSection}{Professional Experience}

% {\bf Postdoctoral Research Assistant, Dr. John Chodera} \hfill {\em September 2021 - June 2022}
% \newline co-designed SPICE dataset (700k drug‑like molecules, 17M conformers), trained transformer FFs achieving less than 1.0 kcal/mol MAE on off‑equilibrium geometries

% {\bf Postdoctoral/Graduate Research Assistant, Dr. Olaf Wiest} \hfill {\em August 2018 - October 2020}
% \newline implemented PyTorch GNN to predict organic reaction products/yields with top‑1 accuracy 18 pp and piloted on pharma proprietary set (>500k rxns), integrated real‑world yield data, delivering 0.27 R² external‑val, code adopted by 2 pharma partners

% {\bf Graduate Research Assistant, Dr. John Parkhill} \hfill {\em August 2015 - August 2018}
% \newline curated 10M‑conformation dataset (45k molecules) for TensorMol NN force‑field, automated conformer sampling \& DFT optimization on 80k‑core cluster, trained NNFF achieving 0.8 kcal/mol energy MAE \& 25 meV/Å force MAE, reduced cpu usage 40\% vs antecedent, co‑authored 7 publications incl. JACS \& Chem. Sci.; first‑author JCP Editor’s Pick on autoencoder chemical embeddings

% \end{rSection}

\begin{rSection}{Professional Experience}

\textbf{Postdoctoral Research Assistant, Dr.\ John Chodera} \hfill \emph{Sept 2021 -- Jun 2022}
\vspace{-0.4em}
\begin{itemize}\itemsep -0.4em
  \item co-built \emph{SPICE} dataset (700 k mols, 17 m confs); trained transformer ff with < 1 kcal/mol off-equilibrium MAE
\end{itemize}

\textbf{Postdoctoral/Graduate Research Assistant, Dr.\ Olaf Wiest} \hfill \emph{Aug 2018 -- Oct 2020}
\vspace{-0.4em}
\begin{itemize}\itemsep -0.4em
  \item built pytorch gnn for rxn outcome/yield → +18 pp top-1 vs baseline
  \item validated on 500 k proprietary rxns (R²=0.27); code adopted by 2 pharma partners
\end{itemize}

\textbf{Graduate Research Assistant, Dr.\ John Parkhill} \hfill \emph{Aug 2015 -- Aug 2018}
\vspace{-0.4em}
\begin{itemize}\itemsep -0.4em
  \item curated 10 m-conf/45 k-mol dataset; automated dft on 80 k-core cluster
  \item trained tensormol nn-ff (0.8 kcal/mol e mae, 25 mev/Å f mae) → –40 % cpu vs prior
  \item led 7 pubs (jacs, chem sci) incl. jcp editor’s pick on chem embeddings
\end{itemize}

\end{rSection}

%----------------------------------------------------------------------------------------
%	Publications
%----------------------------------------------------------------------------------------
\begin{rSection}{Selected Publications}

\item \fullcite{eastman2023spice}
\item \fullcite{saebi2023use}
\item \fullcite{wang2022end}
\item \fullcite{yao2018tensormol}

\end{rSection}

\clearpage

\end{document}